\documentclass{article}

%\include{macros}

\bibliographystyle{plain}

\usepackage{epsfig}
\usepackage{xcolor}
%%\usepackage{amssymb}
\usepackage{amsmath}

\input{abbrev.tex}

\begin{document}

\title{High order finite volume operators in a cut cell context,
  part 2: truncation error measurements}
\author{
  D. T. Graves    \footnotemark[1]
        }

\maketitle

\begin{abstract}

Devendran, et al. \cite{Devendran2017} present an algorithm that
solves Poisson's equation to fourth-order accuracy using an EB grid.
Generalizing somewhat  the framework they present, we can define a
family of elliptic operators.   That algorithm is a special case of a
family of finite volume elliptic operators.   In part 1
\cite{Devendran2023}, we present the formal 
process to generate these operators and present system matrix
conditioning data.      Here we present how truncation error varies
within the family of operators.

\section{Underlying mathematics}
\label{sec::notation}

\subsection{Embedded boundary description}
Embedded boundary (EB) grids are formed when one passes an  surface
through a Cartesian mesh.    For sufficiently complex geometries, these methods
are very attractive because grid generation is a solved problem
even for moving grids \cite{MillerTrebotich2012}.
In the current context, we form the cutting surface as the zero
surface of a function  of space $I(\xbold)$, $\xbold \in R^D$.
For smooth ($I$),  moments can be generated to any
accuracy \cite{Schwartz2015}.

Formally, the underlying description of space
is given by rectangular control volumes on a Cartesian mesh
$\Upsilon_\ibold = [(\ibold-\half {\ubold})h, (\ibold+\half
{\ubold})h], \ibold \in \bigzbold^D$, where $D$ is the dimensionality
of the problem, $h$ is the mesh spacing, and ${\ubold}$ is the vector
whose entries are all one (note we use bold font $\ubold = (u_1, \dots, u_d,
\dots, u_D)$ to indicate a vector quantity).
Given an irregular domain $\Omega$, we
obtain control volumes $V_\ibold = \Upsilon_\ibold \bigcap \Omega$ and
faces $A_{\ibold,d\pm} = A_{\ibold \pm \half \ebold_d}$ which are the
intersection of the boundary of $\partial V_\ibold$ with the
coordinate planes $\{{\xbold}:x_d = (i_d \pm \half)h \}$ ($\ebold_d$ is
the unit vector in the $d$ direction).  We also
define $A_{B,\ibold}$ to be the intersection of the boundary of the
irregular domain with the Cartesian control volume: $A_{B,\ibold}
= \partial \Omega \bigcap \Upsilon_\ibold$. 

\subsection{Notation}

We use the compact notation
\begin{align*}
(\xbold - \xbar)^\pbold &= \prod\limits^D_{d=1} (x_d - {\bar x}_d)^{p_d} \\
\pbold! &= \prod\limits^D_{d = 1} p_d!
\end{align*}
Given a $D$-dimensional region of space $\ibold$, 
and a $D$-dimensional integer vector
$\pbold$, we define $m_\ibold^\pbold(\xbar)$ to be the $\pbold^{th}$
moment of $\ibold$ to the point $\xbar$.
\begin{equation}
\mcal_\ibold^\pbold(\xbar)  =  \int\limits_{\ibold} (\xbold - \xbar)^\pbold dV^D.
\label{eqn::volmoment}
\end{equation}
This breaks out into three different cases, volumes, coordinate faces
and irregular faces.
A coordinate face $f_{12}$  is the area where two
volumes ($V_1, V_2$) intersect.
A irregular face $f_V$  is area where the cutting surface
intersects the volume $V$.
The volume moment for a particular power $\pbold$  is
given by
\begin{equation}
  \mcal^\pbold_V = \int\limits_{\ibold} (\xbold - \xbar)^\pbold dV
  \label{eqn::volmom}
\end{equation}
The $\pbold^{th}$ coordinate face moment for the area of where volumes $V_1, V2$
intersect is given by
\begin{equation}
  \mcal^\pbold_{f_{12}} = \int\limits_{A(V_1\bigcap V_2)} (\xbold - \xbar)^\pbold dA.
  \label{eqn::facmom}
\end{equation}
The cutting surface of the boundary $G$ is given by  the zero
surface of the  implicit function ($I(\xbold)= 0$) used in grid generation.
The $\pbold^{th}$ irregular face moment is given by an integral over
the intersection of this cutting surface and the Euclidean volume:
\begin{equation}
  \mcal^\pbold_V = \int\limits_{A(V \bigcap G)} (\xbold - \xbar)^\pbold dA.
  \label{eqn::irrmom}
\end{equation}
These moments are natural products of the grid generation algorithm
in \cite{Schwartz2015}.


\section{Poisson operator}

Devendran, et al. \cite{Devendran2017} present an algorithm that
solves Poisson's equation to fourth-order accuracy using embedded
boundaries.  Generalizing somewhat  the framework they present, we
present a family of elliptic operators.

Poisson's equation for field $\psi$, given charge $\rho$:
\begin{equation}
  \nabla \cdot (\nabla \psi) = \rho
  \label{eqn::poisson}
\end{equation}
If we integrate equation \ref{eqn::poisson} and apply the divergence theorem 
over control volumes we get
\begin{equation}
\int_{V_\ibold} \nabla \cdot (\nabla \psi ) dV = \int_{\partial
  V_\ibold} \nabla \psi\cdot \nhat  dA = \int_{V_\ibold} \rho dV
\end{equation}
Without approximation, we can say that the surface integral can be
computed by a sum over the faces that comprise the surface of the
volume.
\begin{equation}
\sum_{f: f \in \partial V_\ibold} \nabla \psi \cdot \nhat dA = \int_{V_\ibold} \rho dV
\end{equation}
Using equations \ref{eqn::faceavg} and \ref{eqn::volavg}, the discrete
operator $L_h$ becomes
\begin{equation}
L_h(\phi) = \frac{1}{|V_\ibold|}\sum_{f: f \in \partial V_\ibold} <\nabla \psi \cdot \nhat>_f A_f
\label{eqn::discretePoisson}
\end{equation}
Put in terms of fluxes, 
\begin{equation}
L_h(\phi) = \frac{1}{|V_\ibold|}\sum_{f: f \in \partial V_\ibold} F_f A_f
\end{equation}
where the flux is the integral over the face of the normal gradient
\begin{equation}
  F_f = 
\end{equation }
A finite volume operator is defined by the algorithm to compute the
integral of fluxes at each face.

The truncation error $E_T$ of a discrete  operator $L_h$ when applied to an exact
solution $\phi_e$ is given by
\begin{equation}
  E_T = L_h(\phi_e) - L_e(\phi_e)
\end{equation}
where L_e is exact operator.    how $L_e$ and $\phi_e$
There are several orders of accuracy in play:
\begin{itemize}
\item{P_E} is the order to which the exact solutions are computed.
\item The geometric moments must be computed to $P_E$.
\item The discrete fluxes and so are are computed to $P_C$, the order
  to which the operator is being computed.
\item If we take the obvious choice of $P_E = P_C$, the truncation
  error of the operator will appear too large because the exact
  solution is not smooth enough.  
\item  We will be safe and make $P_E = P_C + 3$.



  Since debugging this test is taking forever and a bunch of things have changed,
  here are a bunch of things I do not want to forget to document some
  of the things I have learned as of 9-15-2025 that I want to be part
  of this document.
  \begin{itemize}
    \item Solving the normal equations is right out.   As the system
      matrix gets super stiff (as is documented in an earlier tech
      report), the accuracy of the solve goes away and you get
      unpredictable answers.   Switching to using SVD (via Eigen) to
      get  a pseudoinverse makes this problem completely go away.   
      Some of the test code has been left in (testing the pseudo
      inverse and so on) has been left in the code.
    \item The operator order $P_O$ and the order of the geometry $P_G$
      can be the same thing but doing this was giving poor answers
      when I compared against exact Taylor coefficients.  Setting
      $P_G = P_O + 2$ makes for much better looking answers.
    \item Moment shifting $(x-{\bar{x}})$  is used in setting up the
      stencil.  For computing the $Q$ matrix, ${\bar{x}}$ does not
      apply.    The moments still may need to be shifted but not by
      the same distance (maybe we need to shift to the face centroid?).
  \end{itemize}



\footnotetext[1]{Lawrence Berkeley National Laboratory, Berkeley,
  CA. Research at LBNL was supported financially by the Office of
  Advanced Scientific Computing Research of the US Department of
  Energy under contract number DE-AC02-05CH11231.}
\renewcommand{\thefootnote}{\fnsymbol{footnote}}\
\bibliography{references}

\end{document}
